\documentclass{dtek}

\usepackage{titlesec}
\titleformat{\paragraph}[hang]{\normalfont\normalsize\bfseries}{\theparagraph}{1em}{}
\titlespacing*{\paragraph}{0pt}{3.25ex plus 1ex minus 0.2ex}{0.7em}

\newcommand{\updated}{2013--03--19}

\title{Ekonomiskt reglemente}
\date{Uppdaterad \updated}

\begin{document}

\makeheadfoot%

\vspace*{\fill}
\begin{center}
  {\Huge \textbf{Ekonomiskt reglemente för Datateknologsektionen}}
  
  \includegraphics[width=300pt]{Datalogo.pdf}
  
  {\LARGE Chalmers, Göteborg}

  \vspace{2cm}

  {\LARGE Uppdaterad \updated}
\end{center}
\vfill

\newpage

\setcounter{tocdepth}{1}
\tableofcontents

\newpage

%Allmänt%%%%%%%%%%%%%%%%%%%%%%%%%%%%%%%%%%%%%%%%%%%%%%%%%%%%%%%%%%%%%%%
\section{Allmänt}
\subsection{}
Detta reglemente har skapats för att skapa tydlighet över vad sektionens pengar får och inte får användas till. Det är därmed tänkt att vara ett hjälpmedel för styrelsen, revisorerna och sektionens föreningar i deras arbete. Grundfilosofin är att sektionens pengar är till för sektionens samtliga medlemmar, vilket i praktiken innebär att arrangemang som inte är öppna för dessa inte får belasta sektionens ekonomi. Viktigt att tänka på vid tolkandet utav reglementet är att det inte har utformats för att i onödan begränsa föreningars möjligheter eller lust till roliga arrangemang och upptåg.
\subsection{Prisbasbelopp}
\begin{itemize}
  \item Det prisbasbelopp som skall användas under hela verksamhetsåret är det prisbasbelopp som är aktuellt vid verksamhetsårets början.
  \item Delar av prisbasbelopp skall avrundas till närmaste hundratal. 
\end{itemize}

%D-Styret%%%%%%%%%%%%%%%%%%%%%%%%%%%%%%%%%%%%%%%%%%%%%%%%%%%%%%%%%%%%%%
\section{D-Styret}
\subsection{}
D-Styret är ansvariga för sektionens hela ekonomi
\subsection{}
Ordförande är skyldig att
\begin{itemize}
  \item Teckna sektionens firma
  \item Se till att ordförande och kassör i varje förening har tillgång till och kunskap om sektionens stadgar, reglemente och ekonomiska reglemente 
\end{itemize}
\subsection{}
Kassören är skyldig att
\begin{itemize}
  \item Teckna sektionens firma
  \item Fortlöpande kontrollera sektionens räkenskaper och bokföring
  \item I samråd med styrelsen upprätta budgetförslag till första ordinarie höstmötet
  \item Till varje sektionsmöte kunna redogöra för sektionens ekonomiska ställning
  \item Utbilda nya sektionsfunktionärer i hur sektionens bokförings och redovisningssystem skall användas
  \item Ansvara för styrelsens handkassa
\end{itemize}
\subsection{}
Styret har rättighet att inom ramen för styrets budget representera sektionen och ordna förtäring till styrets ledamöter i samband med styrelsemöten. Denna förtäring skall vara med som post i styrelsens budget.

%Sektionskommittéer%%%%%%%%%%%%%%%%%%%%%%%%%%%%%%%%%%%%%%%%%%%%%%%%%%%%
\section{Sektionskommittéer}
\label{sec:sektionsforeningar}
\subsection{}
Ordförande i sektionskommitté är skyldig att
\begin{itemize}
\item Kontinuerligt meddela föreningens ekonomiska status till styret
\item Tillsammans med föreningens kassör ansvara för att föreningen förvaltar sina tillgångar i enlighet med sektionens stadgar, reglementen och beslut.
\item Tillsammans med föreningens kassör skriva ett bindande avtal med
styret angående skuldfrågan vid felaktig bokföring
\item Tillsammans med föreningens kassör teckna föreningens firma
\end{itemize}
\subsection{}
Kassör i sektionskommitté är skyldig att
\begin{itemize}
\item Föra kassabok av sådan typ som godkänts av sektionens revisorer
\item Tillsammans med föreningens ordförande ansvara för att föreningen förvaltar sina tillgångar i enlighet med sektionens stadgar, reglementen och beslut.
\item Tillsammans med föreningens ordförande skriva ett förbindande avtal med Styret angående skuldfrågan vid felaktigt förd bokföring
\item På varje sektionsmöte redovisa föreningens ekonomiska situation
\item Arkivera sektionens bokföringar, på en plats anvisad av Styret, så lång tid som föreskrivs för den organisationsform som datateknologsektionen är.
\item Tillsammans med föreningens ordförande teckna föreningens firma
\end{itemize}
\subsection{}
Varje förening skall påbörja en ny kassabok i starten av sitt verksamhetsår
\subsection{Äskning}
Föreningar som önskar extra ekonomiska medel till sin verksamhet skall inkomma med önskemål och skäl därför till styret. Styret kan bevilja extra medel om beloppet understiger 0,25 prisbasbelopp, om beloppet är större måste först sektionsmöte tillfrågas.
\subsection{Kapital}
\label{sec:sektionsforeningar_startkapital}
\subsubsection{Startkapital}
DRUST, DAG, Delta, D6 och DNollK skall vid mandatperiodens början ha 0 prisbasbelopp i tillgångar och ha rätt att från styrelsen låna upp till 0.25 prisbasbelopp.
\subsubsection{Slutkapital}
DRUSTs, DAGs, Deltas, D6s och DNollKs tillgångar samt skulder skall vid slutet av mandatperioden tillfalla styrelsen.
\subsection{Budget}
För arrangemang med omsättning överstigande 0,25 prisbasbelopp skall en noggrann budget upprättas och godkännas av styret innan arrangemanget får genomföras.
\subsection{Arbetskläder}
Föreningar får lägga en rimlig summa pengar på arbetskläder. Rimligt kan anses vara en overall eller motsvarande, t-shirt samt annan tröja/piké.
\subsection{Jobbarmat}
Vid arrangemang inom föreningens verksamhet då det ej finns möjlighet att gå hem och äta får föreningens ekonomi belastas med utgifter för att laga jobbarmat. Dessa utgifter skall vara på en rimlig nivå samt ingå i arrangemangets budget.

%Intresseföreningar och andra grupper under sektionen%%%%%%%%%%%%%%%%%%
\section{Intresseföreningar och andra grupper under sektionen}
\subsection{}
Ekonomiska medel för intresseföreningar och andra grupper under sektionen skall regleras av styrets budget
\subsection{}
I de fall en intresseförening eller annan grupp under sektionen har beviljats egen ekonomi gäller även §\ref{sec:sektionsforeningar}, med undantag för §\ref{sec:sektionsforeningar_startkapital}, för dessa.

%Internrepresentation%%%%%%%%%%%%%%%%%%%%%%%%%%%%%%%%%%%%%%%%%%%%%%%%%%
\section{Internrepresentation}
\subsection{}
\label{sec:internreps}
Varje sektionskommitté har rättighet att varje verksamhetsår bekosta en teambuildingaktivitet och en överlämning med föreningens ekonomi. Dessa bör dock ligga på en rimlig ekonomisk nivå och inte överskrida föreningens ekonomiska resurser. Den totala summan får maximalt uppgå till 0.015 prisbasbelopp per medlem.\\
Alkoholhaltiga drycker får ej bekostas med föreningens pengar vid sådant arrangemang.
\subsection{}
Styrelsen har rätt att bevilja intresseföreningar eller andra grupperingar under teknologsektionen rättighet motsvarande §\ref{sec:internreps} om det anses vara motiverat.

%Sponsring%%%%%%%%%%%%%%%%%%%%%%%%%%%%%%%%%%%%%%%%%%%%%%%%%%%%%%%%%%%%%
\section{Sponsring}
\subsection{}
Föreningar får inte söka sponsring utan samråd med datas
arbetsmarknadsgrupp, Dag.

%Fonder%%%%%%%%%%%%%%%%%%%%%%%%%%%%%%%%%%%%%%%%%%%%%%%%%%%%%%%%%%%%%%%%
\section{Fonder}
\subsection{Kapitalfonden}
\subsubsection{Syfte}
\label{sec:kapitalfond_syfte}
Syftet med kapitalfonden är att avlasta sektionens respektive de olika sektionsföreningarnas ekonomi från stora investeringar. Pengarna skall användas till saker som har ett bestående värde, samt är till gagn för sektionens medlemmar direkt eller indirekt. Pengar skall ej användas till driftbidrag eller stöd för förgänglig verksamhet. Fonden skall ej användas till verksamhet som lokalfond är ämnad för.
\subsubsection{Förvaltning}
Fonden förvaltas av sektionens styrelse.
\subsubsection{Avsättning}
\begin{itemize}
\item En av styrelsen budgeterad summa som godkänts av sektionsmötet
\item All avkastning ifrån kapitalfonden tillförs kapitalfonden.
\end{itemize}
\subsubsection{Uttag}
\begin{itemize}
\item Sektionsstyrelsen har rätt att bevilja uttag ur fonden på belopp upp till totalt 0,25 prisbasbelopp per år. Uttag av belopp överstigande 0,25 prisbasbelopp skall godkännas av sektionsmöte innan medel utbetalas. Varje uttag skall redovisas inför nästkommande sektionsmöte.
\item Sektionsmötet äger rätt att vid välmotiverat behov göra uttag i strid med syftet i §\ref{sec:kapitalfond_syfte}
\end{itemize}
\subsubsection{Revision}
Revidering sker varje år av sektionens revisorer. Revision sammanfaller med övrig på sektionen förekommande.
\subsection{Lokalfonden}
\subsubsection{Syfte}
\label{sec:lokalfond_syfte}
Syftet med lokalfonden är att säkra medel för underhåll och reparationer av sektionens lokaler, samt för inköp av inventarier. Pengarna skall användas till större reparationer och ommålningar av sektionslokalerna, samt möbler till trivselytor där teknologen i gemen har tillträde.
\subsubsection{Förvaltning}
Fonden förvaltas av sektionens styrelse.
\subsubsection{Avsättning}
\label{sec:lokalfond_avsattning}
\begin{itemize}
\item Tio (10) \% av under verksamhetsåret influtna sektionsavgifter tillförs lokalfonden.
\item En av styrelsen budgeterad summa som godkänts av sektionsmötet
\item All avkastning ifrån lokalfonden tillförs lokalfonden.
\end{itemize}
\subsubsection{Uttag}
\begin{itemize}
\item Sektionsstyrelsen disponerar 3/4 av årets tillförda kapital,
enligt första punkten i §\ref{sec:lokalfond_avsattning}, för basdrift av sektionslokalerna. Sektionsstyrelsen har rätt att besluta om uttag på överstigande belopp, dock skall detta redovisas inför nästkommande sektionsmöte. Vid detta sektionsmöte skall i sådana fall även genomförda eller planerade inköp, reparationer och underhåll redovisas.
\item Sektionsmötet äger rätt att vid välmotiverat behov göra uttag i strid med syftet i §\ref{sec:lokalfond_syfte}.
\end{itemize}
\subsubsection{Revision}
Revidering sker varje år av sektionens revisorer. Revision
sammanfaller med övrig på sektionen förekommande.
\end{document}
